% Moderne Farbdefinitionen für LuaLaTeX
% Professionelle Farbpalette für akademische Dokumente

% --- Modernes Farbpaket ---
\usepackage{xcolor}              % Ersetzt veraltetes 'color' Paket

% --- Basis-Farbpalette (Corporate Design) ---
% Hauptfarben
\definecolor{primary}{RGB}{0, 102, 153}        % Professionelles Blau
\definecolor{secondary}{RGB}{102, 102, 102}    % Neutrales Grau
\definecolor{accent}{RGB}{204, 51, 0}          % Akzentfarbe (Orange-Rot)

% --- Erweiterte Grautöne (für professionelle Layouts) ---
\definecolor{darkgray}{RGB}{64, 64, 64}        % Fast Schwarz
\definecolor{mediumgray}{RGB}{128, 128, 128}   % Mittelgrau
\definecolor{lightgray}{RGB}{192, 192, 192}    % Hellgrau
\definecolor{verylightgray}{RGB}{240, 240, 240} % Sehr hell (Hintergründe)

% --- Semantische Farben ---
\definecolor{success}{RGB}{40, 167, 69}        % Grün (Erfolg)
\definecolor{warning}{RGB}{255, 193, 7}        % Gelb (Warnung)
\definecolor{danger}{RGB}{220, 53, 69}         % Rot (Gefahr)
\definecolor{info}{RGB}{23, 162, 184}          % Cyan (Information)

% --- Code-Highlighting Farben (abgestimmt mit listings) ---
\definecolor{codebackground}{RGB}{248, 248, 248}
\definecolor{codekeyword}{RGB}{0, 102, 153}    % = primary
\definecolor{codecomment}{RGB}{102, 102, 102}  % = secondary
\definecolor{codestring}{RGB}{153, 153, 0}
\definecolor{codenumber}{RGB}{153, 0, 85}

% --- Druck-optimierte Farben ---
% CMYK für professionellen Druck
\definecolor{printblack}{cmyk}{0,0,0,0.85}     % Lesbareres Schwarz
\definecolor{printgray}{cmyk}{0,0,0,0.30}      % Mittleres Grau
\definecolor{printblue}{cmyk}{1,0.5,0,0}       % Druckbares Blau

% --- Universitäts-/Corporate-Farben (anpassbar) ---
% Beispiel für Universität der Bundeswehr München
\definecolor{unibwblue}{RGB}{0, 51, 102}       % UniBw Blau
\definecolor{unibwgray}{RGB}{102, 102, 102}    % UniBw Grau
% Passen Sie diese an Ihre Institution an!

% --- Barrierefreie Farben (WCAG 2.1 AA konform) ---
\definecolor{accessible-blue}{RGB}{0, 73, 144}
\definecolor{accessible-green}{RGB}{0, 128, 55}
\definecolor{accessible-red}{RGB}{187, 25, 25}
\definecolor{accessible-orange}{RGB}{255, 140, 0}

% --- Moderne Farbverläufe (für TikZ/Grafiken) ---
\definecolor{gradient-start}{RGB}{0, 102, 153}
\definecolor{gradient-end}{RGB}{102, 204, 255}

% --- Konditionale Farbwahl: Druck vs. Digital ---
\newif\ifprintversion
% \printversiontrue   % Für Druck
\printversionfalse    % Für Digital (Standard)

\ifprintversion
  % === DRUCK-VERSION ===
  \colorlet{textcolor}{printblack}
  \colorlet{linkcolor}{printblue}
  \colorlet{headingcolor}{printblack}
  \color{textcolor}                % Setze Standardtextfarbe
\else
  % === DIGITAL-VERSION ===
  \colorlet{textcolor}{black}
  \colorlet{linkcolor}{primary}
  \colorlet{headingcolor}{primary}
  % Standardfarbe bleibt schwarz für bessere Bildschirmlesbarkeit
\fi

% --- Utility-Befehle ---
% Einfache Farbwechsel
\newcommand{\textprimary}[1]{\textcolor{primary}{#1}}
\newcommand{\textsecondary}[1]{\textcolor{secondary}{#1}}
\newcommand{\textaccent}[1]{\textcolor{accent}{#1}}
\newcommand{\textwarning}[1]{\textcolor{warning}{#1}}
\newcommand{\textdanger}[1]{\textcolor{danger}{#1}}
\newcommand{\textsuccess}[1]{\textcolor{success}{#1}}

% Hintergrundfarben
\newcommand{\highlight}[1]{\colorbox{verylightgray}{#1}}
\newcommand{\codehighlight}[1]{\colorbox{codebackground}{\texttt{#1}}}

% --- KOMA-Script Integration ---
% Überschriften einfärben
\addtokomafont{chapter}{\color{headingcolor}}
\addtokomafont{section}{\color{headingcolor}}
\addtokomafont{subsection}{\color{headingcolor}}

% --- TikZ-Integration ---
\usepackage{tikz}
% Farbstile für TikZ definieren
\tikzset{
  primary/.style={color=primary, fill=primary!10},
  secondary/.style={color=secondary, fill=secondary!10},
  accent/.style={color=accent, fill=accent!10},
  success/.style={color=success, fill=success!10},
  warning/.style={color=warning, fill=warning!10},
  danger/.style={color=danger, fill=danger!10}
}

% --- Debug-Information ---
\makeatletter
\AtBeginDocument{%
  \typeout{=== Farbpalette geladen ===}
  \ifprintversion
    \typeout{Modus: DRUCK (CMYK-optimiert)}
    \typeout{Textfarbe: printblack}
  \else
    \typeout{Modus: DIGITAL (RGB-optimiert)}
    \typeout{Textfarbe: standard schwarz}
  \fi
  \typeout{Primärfarbe: RGB(0,102,153)}
  \typeout{Akzentfarbe: RGB(204,51,0)}
  \typeout{=============================}
}
\makeatother

% --- Farbpaletten-Showcase (für Tests) ---
% \newcommand{\showcolorpalette}{%
%   \begin{center}
%   \begin{tabular}{|l|c|c|}
%   \hline
%   \textbf{Farbe} & \textbf{RGB} & \textbf{Beispiel} \\
%   \hline
%   Primary & (0,102,153) & \textcolor{primary}{\rule{2cm}{0.5cm}} \\
%   Secondary & (102,102,102) & \textcolor{secondary}{\rule{2cm}{0.5cm}} \\
%   Accent & (204,51,0) & \textcolor{accent}{\rule{2cm}{0.5cm}} \\
%   Success & (40,167,69) & \textcolor{success}{\rule{2cm}{0.5cm}} \\
%   Warning & (255,193,7) & \textcolor{warning}{\rule{2cm}{0.5cm}} \\
%   Danger & (220,53,69) & \textcolor{danger}{\rule{2cm}{0.5cm}} \\
%   \hline
%   \end{tabular}
%   \end{center}
% }